\documentclass[letterpaper]{article}
\usepackage[letterpaper,margin=1in]{geometry}
\input{wsp/testing/test-case-library.tex}
\begin{document}
\def\TCID{TC-0020}
\def\TCTitle{Interactive foreground interruption and cleanup}
\def\TCPurpose{Verify a real console control event cancels and collects the foreground child tree before returning to the prompt.}
\def\TCPriority{\TCPriorityCritical}
\def\TCRequirementRef{WSH-REQ-0020, WSH-REQ-0024, WSH-REQ-0050, WSH-REQ-0056}
\def\TCDesignRef{ADR-0010 and WSH-SPEC-INTERACTIVE-0001 sections 6--7}
\def\TCFeatureRef{Ctrl+C, Ctrl+Break, foreground groups, sigint, prompt recovery}
\def\TCTestTechnique{State-transition and native console integration testing.}
\def\TCPreconditions{The native console driver and runtime probe are built.}
\def\TCEnvironment{Native current Windows; legacy representatives remain a release matrix gate.}
\def\TCAssumptions{The operating system delivers console control events to the shared process group.}
\def\TCInputData{Long-running foreground child, pipeline, Ctrl+C, and Ctrl+Break.}
\def\TCInitialState{No tracked foreground or background child exists.}
\def\TCProcedure{Launch WSH on a genuine console, start work, deliver each control event, and inspect status, prompt recovery, handler order, and process liveness.}
\def\TCExpectedResult{Every live stage receives status 130, no child remains, sigint runs after cleanup, and a fresh prompt accepts input.}
\def\TCPostConditions{The console mode, handler, and all child handles are restored or released.}
\TCTemplate
\end{document}
